\chapter{Lobectomy}

\section*{What Is This Surgery?}
A lobectomy is a major surgical procedure involving the complete anatomical removal of one of the five lobes of the human lung. This intricate resection is most frequently performed for the definitive treatment of early-stage non-small cell lung cancer (NSCLC) that is localized to a single lobe, aiming for curative intent. The procedure's primary goal is to achieve complete tumor excision with clear surgical margins while meticulously preserving as much healthy lung tissue as possible to maintain adequate postoperative respiratory function. It represents a cornerstone in thoracic oncology, offering the best chance for cure in appropriately selected patients. Beyond malignancy, lobectomy is also indicated for certain benign conditions that are refractory to medical management or pose significant risk, such as severe localized infections, extensive bronchiectasis, or large pulmonary arteriovenous malformations.

\section*{Alternative Names}
\begin{itemize}
  \item Pulmonary Lobectomy
  \item Lung Lobe Resection
  \item Anatomical Lobectomy
  \item Video-Assisted Thoracoscopic Surgery (VATS) Lobectomy
  \item Robotic-Assisted Thoracoscopic Surgery (RATS) Lobectomy
  \item Open Thoracotomy Lobectomy
\end{itemize}

\section*{Relevant Surgical Anatomy}
The human lungs are paired organs situated within the thoracic cavity, separated by the mediastinum. The right lung has three lobes (upper, middle, and lower), while the left lung has two (upper and lower), with the lingula being a functional equivalent of the right middle lobe. Each lobe is a distinct anatomical and functional unit, supplied by its own segmental bronchus, pulmonary artery branch, and pulmonary vein branch, which collectively form the hilum of the lung. The lobes are separated by fissures: the oblique fissure separates the upper and middle lobes from the lower lobe on the right, and the upper lobe from the lower lobe on the left; the horizontal fissure on the right separates the upper from the middle lobe. These fissures can be incomplete, requiring surgical stapling for full separation.

Key structures within the hilum, which must be meticulously identified and ligated during lobectomy, include the main pulmonary artery and its lobar branches, the superior and inferior pulmonary veins and their lobar tributaries, and the main bronchus and its lobar divisions. The lymphatic drainage of the lungs follows the bronchial tree, draining into intrapulmonary, hilar (peribronchial), and then mediastinal lymph nodes (e.g., subcarinal, paratracheal, aortopulmonary window nodes). The phrenic nerve runs anterior to the lung hilum, innervating the diaphragm, while the vagus nerve runs posterior. The recurrent laryngeal nerve, a branch of the vagus, loops around the aorta on the left and the subclavian artery on the right, ascending to the larynx, and is vulnerable to injury during mediastinal lymph node dissection, particularly in the aortopulmonary window.

\section{Surgical Principle \& Rationale}
The fundamental principle of a lobectomy lies in the anatomical segmentation of the lungs, where each lobe functions as a distinct unit, supplied by its own segmental bronchus, pulmonary artery branch, and pulmonary vein branch. The surgical goal is to remove the diseased lobe en bloc, along with its associated hilar lymph nodes and its dedicated vascular and bronchial supply, ensuring complete resection of the pathology. This approach, known as an anatomical resection, is crucial for oncological cases to achieve negative surgical margins (R0 resection) and facilitate comprehensive systematic mediastinal lymph node dissection, which is vital for accurate staging, prognosis, and local disease control.

The rationale behind lobectomy is to balance radical disease removal with the preservation of functional lung parenchyma. By individually ligating and dividing only the specific hilar structures (bronchus, artery, vein) supplying the affected lobe, the remaining healthy lobes can re-expand and continue to function, minimizing the impact on overall respiratory capacity. This selective removal distinguishes lobectomy from less extensive resections like wedge resections or segmentectomies, which may be appropriate for smaller, peripheral lesions or patients with severely compromised lung function, and from more extensive resections like pneumonectomy, which involves removing an entire lung and carries higher morbidity.

\section{History \& Surgical Evolution}
The evolution of pulmonary lobectomy is a testament to advancements in surgical technique, anesthesia, and understanding of thoracic anatomy. Early attempts at lung resection in the late 19th and early 20th centuries were fraught with extremely high mortality rates due to uncontrolled hemorrhage, infection, and respiratory failure. The first successful one-stage pneumonectomy for lung cancer, a more extensive procedure than lobectomy, is often attributed to Dr. Evarts Graham in 1933, which paved the way for safer anatomical resections by demonstrating the feasibility of major lung surgery. Subsequent pioneers like Dr. Richard Overholt and Dr. Edward Churchill refined lobectomy techniques in the 1940s and 1950s, standardizing the individual ligation of hilar structures and significantly improving postoperative care.

The advent of modern general anesthesia, particularly the development of one-lung ventilation techniques using double-lumen endotracheal tubes, significantly reduced intraoperative risks and improved surgical exposure. The most transformative development in recent decades has been the introduction of minimally invasive approaches. Video-assisted thoracoscopic surgery (VATS) lobectomy was pioneered in the early 1990s, notably by Dr. Robert J. McKenna Jr. and Dr. Rodney J. Landreneau, offering reduced pain, shorter hospital stays, and faster recovery compared to traditional open thoracotomy, without compromising oncological outcomes. Robotic-assisted thoracoscopic surgery (RATS) further refines this minimally invasive paradigm, providing enhanced dexterity, tremor filtration, and 3D visualization, particularly for complex hilar dissections.

\section{Clinical Indications}
Lobectomy is indicated for a range of pulmonary pathologies, primarily malignant but also for select benign conditions that are localized to a single lobe. The decision to proceed with lobectomy is based on careful evaluation of the patient's overall health, pulmonary function, and the nature and extent of the disease.

\begin{itemize}
  \item **Early-stage Non-Small Cell Lung Cancer (NSCLC):** This is the most common indication, particularly for Stage I and II tumors confined to a single lobe, where complete surgical resection offers the best chance for cure.
  \item **Carcinoid Tumors:** Resection of typical or atypical carcinoid tumors, which are neuroendocrine tumors of the lung, often requires lobectomy to achieve clear margins due to their potential for local invasion.
  \item **Metastatic Tumors to the Lung:** In highly selected patients with oligometastatic disease from extrapulmonary primary cancers, lobectomy may be performed to resect solitary or few metastases, especially if the primary tumor is controlled and the patient has adequate pulmonary reserve.
  \item **Severe Localized Bronchiectasis:** For patients with chronic, irreversible dilation and infection of the bronchi in a specific lobe, refractory to medical therapy, lobectomy can alleviate symptoms and prevent recurrent infections.
  \item **Chronic Lung Abscess or Necrotizing Pneumonia:** In rare cases where a severe, localized lung infection or abscess does not respond to prolonged antibiotic therapy and poses a life-threatening risk, surgical debridement via lobectomy may be necessary.
  \item **Pulmonary Arteriovenous Malformations (AVMs):** Large or symptomatic AVMs, which can cause hypoxemia, hemoptysis, or paradoxical emboli, may necessitate lobectomy for definitive treatment.
  \item **Fungal Infections (e.g., Aspergilloma):** Localized fungal balls (aspergillomas) within lung cavities, especially if causing recurrent or severe hemoptysis, can be an indication for lobectomy.
  \item **Benign Tumors:** Large or symptomatic benign tumors (e.g., hamartomas, pleomorphic adenomas) that cannot be managed by less extensive resections due to size or central location.
\end{itemize}

\section*{Clinical Positioning}
Lobectomy is overwhelmingly considered a primary, definitive treatment for early-stage non-small cell lung cancer and many other localized pulmonary pathologies. For NSCLC, it is the gold standard surgical intervention, aiming for curative resection, particularly for Stage I and II disease. In this context, it is typically the first-line treatment offered to patients who are medically fit and whose disease is amenable to complete surgical removal, often following multidisciplinary tumor board discussion.

However, lobectomy can also play a secondary or adjunctive role in specific clinical scenarios. For instance, in patients with locally advanced lung cancer (Stage III) who have undergone neoadjuvant chemotherapy or chemoradiation to downstage the tumor, lobectomy may be performed as a secondary procedure to achieve complete pathological response or R0 resection. Similarly, in cases of recurrent benign disease or complications following less extensive resections (e.g., persistent air leak after segmentectomy), a lobectomy might be considered a secondary, salvage procedure. For metastatic disease, it is often secondary to systemic therapy and reserved for highly selected patients with limited metastatic burden and controlled primary disease, forming part of a multimodal treatment strategy.

\section*{Surgical Technique \& Procedure}
Lobectomy is a complex surgical procedure requiring meticulous technique and a thorough understanding of thoracic anatomy. It can be performed via an open thoracotomy or, increasingly, through minimally invasive approaches such as VATS or RATS. The choice of approach depends on tumor characteristics, patient comorbidities, and surgeon expertise.

The procedure begins with the patient undergoing general anesthesia, typically with a double-lumen endotracheal tube or bronchial blocker to facilitate one-lung ventilation. This allows the surgical lung to be deflated, improving exposure and access to the hilar structures. The patient is usually positioned in a lateral decubitus position, with the operative side facing upwards, and the arm abducted and flexed.

The key operative steps generally include:
\begin{enumerate}
  \item **Access and Exploration:** For an open thoracotomy, a posterolateral incision is made, typically extending from the mid-axillary line to the posterior scapular line, allowing for rib spreading and direct visualization. For VATS or RATS, several small incisions (ports) are made, usually 2-4, through which a camera and specialized instruments are inserted. The pleural cavity is entered, and the lung and surrounding structures are inspected to confirm resectability and rule out unforeseen disease spread. Adhesions are carefully lysed.
  \item **Hilar Dissection:** The pulmonary ligament is divided to mobilize the lower lobe. The pleura overlying the hilum is incised, and the individual hilar structures (pulmonary artery branches, pulmonary vein branches, and bronchus) supplying the target lobe are meticulously identified and dissected free from surrounding tissues.
  \item **Ligation and Division of Structures:** Each identified hilar structure is individually ligated with sutures or clips and then divided using surgical staplers or ligatures. The order of division may vary based on surgeon preference and anatomical considerations, but typically the pulmonary vein is addressed first, followed by the artery, and then the bronchus.
  \item **Fissure Completion:** The interlobar fissures, which separate the lobes, are often incomplete and require stapling or division with energy devices to fully separate the diseased lobe from adjacent healthy lobes.
  \item **Lobe Removal:** Once all hilar structures are divided and the lobe is completely separated, it is placed into an endoscopic retrieval bag and removed through one of the incisions (or the main thoracotomy incision).
  \item **Lymph Node Dissection:** A systematic mediastinal lymph node dissection (SMLND) is performed for oncological cases, removing lymph nodes from specific stations (e.g., N2 stations for NSCLC) to ensure accurate staging and improve local control.
  \item **Hemostasis and Air Leak Check:** The surgical field is meticulously inspected for any bleeding points, which are coagulated. The remaining lung is re-inflated under saline to check for any air leaks from the staple lines or lung parenchyma.
  \item **Closure:** After ensuring hemostasis and checking for air leaks, one or two chest tubes are placed into the pleural cavity to drain fluid and air and facilitate re-expansion of the remaining lung. The incisions are then closed in layers using absorbable sutures for muscle and fascia, and skin staples or sutures for the epidermis.
\end{enumerate}

\section*{Preoperative Workup \& Preparation}
Thorough preoperative evaluation is paramount to optimize patient outcomes and minimize risks associated with lobectomy. This comprehensive workup assesses the patient's overall health, cardiac and pulmonary reserve, and the extent of the pulmonary pathology.

Key preparatory steps include:
\begin{itemize}
  \item **Comprehensive Medical History and Physical Examination:** To identify comorbidities, assess functional status, and evaluate for any signs of infection or systemic illness.
  \item **Pulmonary Function Testing (PFTs):** Spirometry, lung volumes, and diffusion capacity (DLCO) are essential to predict postoperative lung function and ensure the patient can tolerate the loss of a lobe. A predicted postoperative forced expiratory volume in 1 second (FEV1) and DLCO of greater than 40\% are generally considered acceptable for lobectomy.
  \item **Cardiovascular Assessment:** Electrocardiogram (ECG), echocardiogram, and sometimes stress testing or cardiac catheterization are performed to evaluate cardiac risk, especially in older patients or those with a history of heart disease.
  \item **Imaging Studies:**
    \begin{itemize}
      \item **Computed Tomography (CT) Scan of the Chest with Contrast:** To precisely localize the lesion, assess its relationship to hilar structures, and evaluate for mediastinal lymphadenopathy or pleural involvement.
      \item **Positron Emission Tomography (PET-CT) Scan:** For oncological cases, to detect distant metastases and assess mediastinal lymph node involvement, guiding staging and treatment decisions.
      \item **Brain MRI:** May be performed for NSCLC staging to rule out brain metastases, particularly for advanced stages or specific histological subtypes.
    \end{itemize}
  \item **Mediastinal Staging:** For lung cancer, invasive staging procedures like endobronchial ultrasound (EBUS) with fine-needle aspiration (FNA) or mediastinoscopy may be required to sample suspicious mediastinal lymph nodes and confirm N2 disease, which might alter treatment strategy (e.g., neoadjuvant therapy).
  \item **Smoking Cessation:** Patients who smoke are strongly advised to quit several weeks prior to surgery to improve pulmonary function, reduce postoperative complications (e.g., pneumonia, wound healing issues), and enhance long-term outcomes.
  \item **Nutritional Optimization:** Addressing malnutrition can improve wound healing, immune function, and reduce infection risk.
  \item **Medication Review:** Anticoagulants and antiplatelet agents are typically held for a specified period before surgery according to institutional protocols. Other medications are reviewed and adjusted as needed.
  \item **Preoperative Education and Consent:** Detailed discussion of the procedure, potential risks, benefits, and alternatives, followed by informed consent, ensuring the patient fully understands the implications.
  \item **Preoperative Antibiotics:** Administered intravenously within 60 minutes prior to incision to reduce the risk of surgical site infection.
\end{itemize}

\section*{Contraindications \& Precautions}
While lobectomy is a highly effective procedure, certain conditions may preclude its safe performance or necessitate extreme caution. These are generally categorized into absolute and relative contraindications.

\textbf{Absolute Contraindications:}
\begin{itemize}
  \item **Inadequate Cardiopulmonary Reserve:** Predicted postoperative FEV1 or DLCO less than 30-40\% of normal, or severe cardiac disease (e.g., recent myocardial infarction, unstable angina, severe pulmonary hypertension) that renders the patient unable to tolerate one-lung ventilation or the physiological stress of surgery.
  \item **Extensive Metastatic Disease:** Evidence of widespread distant metastases (e.g., brain, bone, liver, adrenal glands) that makes curative resection impossible and where surgery would not offer significant palliative benefit.
  \item **Nonsurgical Disease:** Pathologies that are better managed with non-surgical therapies (e.g., extensive small cell lung cancer, diffuse interstitial lung disease, or conditions with widespread bilateral involvement).
  \item **Uncontrolled Bleeding Diathesis:** Severe coagulopathy that cannot be corrected preoperatively, posing an unacceptable risk of hemorrhage.
  \item **Patient Refusal:** The patient declines the procedure after thorough informed consent.
\end{itemize}

\textbf{Relative Contraindications and Precautions:}
\begin{itemize}
  \item **Advanced Age and Frailty:** While not an absolute contraindication, older, frail patients may have increased surgical risks and require more intensive preoperative optimization and careful risk-benefit assessment.
  \item **Extensive Pleural Adhesions:** Previous thoracic surgery, pleurodesis, or inflammatory conditions can lead to dense adhesions, making dissection challenging and increasing the risk of bleeding or lung injury.
  \item **Hilar or Mediastinal Involvement:** Tumors invading critical structures like the main pulmonary artery, superior vena cava, aorta, or carina may require more complex resections (e.g., sleeve lobectomy, pneumonectomy) or may be deemed unresectable.
  \item **Active Infection:** While lobectomy can treat chronic infections, acute, uncontrolled systemic infections or severe pneumonia in the remaining lung can increase surgical risk and should ideally be treated before surgery.
  \item **Immunosuppression:** Patients on immunosuppressive therapy may have a higher risk of infection and impaired wound healing, requiring careful management and potential adjustment of medications.
  \item **Severe Obesity:** Can complicate surgical access, anesthesia management, and increase the risk of respiratory complications such as atelectasis and pneumonia.
  \item **Prior Radiation Therapy to the Chest:** Can cause tissue fibrosis, making dissection more difficult and increasing the risk of complications like bronchopleural fistula due to impaired healing.
  \item **Poor Nutritional Status:** Malnutrition can impair wound healing and increase susceptibility to infection.
\end{itemize}

\section*{Complications \& Risks}
Lobectomy, like any major surgical procedure, carries inherent risks, some of which are general to all surgeries and others specific to thoracic surgery. Patients are thoroughly counseled on these potential complications during the informed consent process.

\begin{itemize}
  \item **General Surgical Risks:**
    \begin{itemize}
      \item **Bleeding (Hemorrhage):** Can occur intraoperatively or postoperatively, potentially requiring blood transfusions or reoperation. Injury to major hilar vessels is a significant concern.
      \item **Infection:** Surgical site infection, pneumonia, or empyema (pus in the pleural space) can occur, sometimes requiring prolonged antibiotic therapy or further surgical drainage.
      \item **Anesthesia Complications:** Risks associated with general anesthesia include adverse reactions to medications, respiratory depression, cardiac events (e.g., myocardial infarction, stroke), and nerve damage from positioning.
      \item **Deep Vein Thrombosis (DVT) and Pulmonary Embolism (PE):** Blood clots can form in the legs and travel to the lungs, a potentially life-threatening complication. Prophylaxis with anticoagulants and mechanical compression devices is routinely used.
    \end{itemize}
  \item **Procedure-Specific Risks:**
    \begin{itemize}
      \item **Prolonged Air Leak:** The most common complication, where air continues to leak from the lung parenchyma through the chest tube for more than 5-7 days. This can prolong hospital stay and occasionally requires further intervention.
      \item **Bronchopleural Fistula (BPF):** A rare but severe complication where a communication forms between the bronchial stump and the pleural space, leading to persistent air leak, empyema, and potential respiratory failure. It often requires reoperation.
      \item **Cardiac Arrhythmias:** Atrial fibrillation is common post-thoracotomy, particularly in older patients, and usually responds to medication, but can prolong hospital stay.
      \item **Pneumonia/Atelectasis:** Collapse of remaining lung tissue or infection can occur due to impaired cough, pain, or retained secretions, necessitating aggressive pulmonary hygiene.
      \item **Respiratory Failure:** In patients with marginal lung function, the loss of a lobe can lead to acute respiratory insufficiency requiring ventilatory support, sometimes prolonged.
      \item **Chylothorax:** Leakage of lymphatic fluid (chyle) into the pleural space due to injury to the thoracic duct, requiring drainage, dietary modification, and sometimes surgical intervention.
      \item **Recurrent Laryngeal Nerve Injury:** Can cause vocal cord paralysis and hoarseness, potentially leading to aspiration and dysphagia.
      \item **Phrenic Nerve Injury:** Can cause diaphragmatic paralysis, leading to shortness of breath and reduced lung capacity.
      \item **Post-thoracotomy Pain Syndrome (PTPS):** Chronic neuropathic pain along the incision site, which can persist for months or years, requiring specialized pain management.
      \item **Incomplete Resection/Positive Margins:** In oncological cases, if tumor cells are found at the edge of the resected tissue, indicating that some cancer may have been left behind, potentially requiring adjuvant therapy or re-resection.
    \end{itemize}
\end{itemize}

\section*{Postoperative Recovery Timeline}
Immediately following a lobectomy, the patient is transferred to the Post-Anesthesia Care Unit (PACU) for close monitoring of vital signs, pain levels, and chest tube output. Pain management is a critical component of immediate postoperative care, often involving epidural analgesia, patient-controlled analgesia (PCA), or multimodal oral pain regimens to facilitate deep breathing and early mobilization. Chest tubes are typically in place to drain fluid and air from the pleural space, allowing the remaining lung to re-expand. Early mobilization, usually within 24 hours, is strongly encouraged to prevent complications such as deep vein thrombosis, pneumonia, and atelectasis. Pulmonary hygiene, including incentive spirometry and coughing exercises, begins immediately.

Over the first 24-48 hours, the focus remains on optimizing pain control, aggressive pulmonary hygiene, and monitoring for complications like excessive bleeding, persistent air leaks, or cardiac arrhythmias. Chest tube output and air leak status are closely monitored; the chest tubes are typically removed when the air leak has resolved, and fluid drainage is minimal (usually <200 mL/24 hours). The hospital stay for a lobectomy usually ranges from 3 to 7 days, depending on the patient's recovery trajectory and the resolution of chest tube issues. Upon discharge, patients receive detailed instructions on wound care, pain management, activity restrictions (e.g., no heavy lifting for 4-6 weeks), and warning signs to report to their surgical team. Long-term recovery involves a gradual return to normal activities over several weeks to months, often supported by pulmonary rehabilitation to improve lung function and exercise tolerance.

\section*{Surgical Findings \& Pathology}
During a lobectomy, the surgeon meticulously documents intraoperative findings, including the precise location and size of the lesion, its relationship to surrounding structures, the presence of any pleural adhesions, and the appearance of hilar and mediastinal lymph nodes. Any unexpected findings, such as additional nodules or evidence of more extensive disease, are also noted. The resected lung lobe and all dissected lymph nodes are then sent for pathological examination.

The pathology report is the most crucial result, particularly for oncological indications. It provides definitive information about the type of tumor (e.g., adenocarcinoma, squamous cell carcinoma), its histological grade, size, and the presence or absence of invasion into surrounding structures (e.g., visceral pleura, chest wall). Critically, it confirms the status of the surgical margins (negative margins, or R0 resection, indicate no cancer cells at the edge of the resected tissue, suggesting complete removal) and the involvement of lymph nodes (N-stage), which is essential for accurate pathological staging (pTNM). For benign conditions, the pathology report confirms the diagnosis (e.g., bronchiectasis, fungal infection, hamartoma) and rules out malignancy. These findings are interpreted by a multidisciplinary team to guide further treatment decisions and determine prognosis.

\section{Expected Postoperative Course}
A normal, successful postoperative course following lobectomy is characterized by a progressive and relatively smooth recovery, indicating effective treatment and an absence of major complications. Patients can expect their pain to be well-controlled with oral analgesics within a few days, allowing for comfortable deep breathing and mobilization. The chest tubes are typically removed within 3-7 days as air leaks resolve and fluid drainage diminishes, leading to full re-expansion of the remaining lung lobes, confirmed by chest X-ray.

Patients should experience a gradual improvement in shortness of breath and fatigue over several weeks, as their body adapts to the reduced lung volume and the surgical trauma heals. Wound healing progresses without signs of infection, and the incision site becomes less tender. Most patients are able to resume light daily activities within 2-3 weeks and gradually return to more strenuous activities, including work, within 6-12 weeks, often with the aid of pulmonary rehabilitation. A successful course means the primary pathology has been definitively addressed, and the patient is on a clear path to regaining their baseline functional status with minimal long-term sequelae.

\section{Postoperative Care \& Follow-up}
Post-lobectomy follow-up is crucial for monitoring recovery, detecting potential complications, and, for oncological patients, surveillance for recurrence. The follow-up schedule is tailored to individual patient needs and the underlying pathology.

\begin{itemize}
  \item **Initial Postoperative Visits:** Typically scheduled 1-2 weeks after discharge for wound check, chest X-ray, and removal of any remaining sutures or staples.
  \item **Pulmonary Rehabilitation:** Many patients benefit from a structured pulmonary rehabilitation program to improve lung function, exercise tolerance, and overall quality of life, especially those with pre-existing lung disease.
  \item **Long-term Oncological Surveillance:** For lung cancer patients, regular follow-up imaging (CT scans of the chest and abdomen) is performed at prescribed intervals (e.g., every 6 months for the first 2-3 years, then annually for up to 5 years) to monitor for local recurrence or distant metastases. PET-CT or brain MRI may be used as indicated.
  \item **Pulmonary Function Testing:** Repeat PFTs may be performed several months post-surgery to assess the impact of the lobectomy on lung function and guide rehabilitation efforts.
  \item **Activity Resumption:** Gradual increase in physical activity is encouraged, with specific guidance on when to resume strenuous activities, heavy lifting, and driving.
  \item **Nutritional Support:** Ongoing dietary advice to support healing and maintain overall health, particularly for patients who experienced significant weight loss.
  \item **Warning Signs:** Patients are educated on warning signs that require immediate medical attention, such as increasing shortness of breath, fever, worsening chest pain, significant wound drainage, hemoptysis, or new neurological symptoms.
\end{itemize}
Regular communication with the surgical and oncology teams is vital for optimal long-term management and addressing any concerns promptly.

\section*{Epidemiology}
Lobectomy is one of the most frequently performed major thoracic surgical procedures worldwide, primarily driven by the high incidence of lung cancer. Lung cancer remains the leading cause of cancer-related mortality globally, with non-small cell lung cancer (NSCLC) accounting for approximately 85\% of all cases. A significant proportion of these cases, particularly those diagnosed at early stages (Stage I and II), are amenable to curative surgical resection via lobectomy. The incidence of lung cancer is higher in older individuals, with the median age at diagnosis typically in the late 60s or early 70s, and it affects both men and women, though historically more prevalent in men due to smoking patterns. The increasing adoption of lung cancer screening programs (low-dose CT scans) in high-risk populations has led to an earlier diagnosis of smaller, resectable tumors, further increasing the number of lobectomies performed. While the overall perioperative mortality rate for lobectomy has significantly decreased over the past few decades due to improved surgical techniques, anesthesia, and postoperative care, it still carries a perioperative mortality rate typically ranging from 1-3\% in high-volume centers, with morbidity rates around 20-40\%. The most common indications for lobectomy remain primary lung cancer, followed by metastatic disease, and then benign conditions such as severe bronchiectasis or fungal infections.

\section*{Outcomes \& Prognosis}
The outcomes and prognosis following lobectomy are highly dependent on the underlying indication, the pathological stage of disease (for cancer), and the patient's overall health and comorbidities. For early-stage non-small cell lung cancer (NSCLC), lobectomy offers the best chance for long-term survival and cure. Five-year survival rates for Stage I NSCLC treated with lobectomy range from 60-80\%, and for Stage II, they are typically 40-50\%. Key prognostic factors for cancer patients include the pathological stage (especially lymph node involvement), tumor size, histological subtype, and the presence of positive surgical margins (R1 or R2 resection). Recurrence rates vary but are generally lower for earlier stages, with most recurrences occurring within the first 2-3 years post-surgery. Functional outcomes are generally good, with most patients experiencing a return to near-baseline activity levels, though some degree of dyspnea on exertion may persist, particularly in those with pre-existing lung disease. Landmark trials like the Lung Cancer Study Group (LCSG) 821 established lobectomy as superior to sublobar resection for early-stage NSCLC. Minimally invasive approaches (VATS/RATS) have demonstrated comparable oncological outcomes to open thoracotomy but with improved short-term recovery metrics, including reduced pain, shorter hospital stays, and faster return to daily activities. For benign conditions, lobectomy typically provides excellent long-term relief of symptoms and resolution of the underlying pathology, with a very good prognosis for functional recovery.

\section{References}
\begin{enumerate}
  \item MedlinePlus. Lung lobectomy. U.S. National Library of Medicine; 2024. Available from: \url{https://medlineplus.gov/ency/article/002996.htm}
  \item Detterbeck FC, et al. Lung Cancer. In: Townsend CM Jr, et al., editors. Sabiston Textbook of Surgery. 21st ed. Elsevier; 2021. p. 1827-1866.
  \item National Comprehensive Cancer Network (NCCN). NCCN Clinical Practice Guidelines in Oncology: Non-Small Cell Lung Cancer. Version 4.2024. Available from: \url{https://www.nccn.org/guidelines/category_1}
  \item Society of Thoracic Surgeons (STS). STS Clinical Practice Guidelines for the Management of Lung Cancer. Available from: \url{https://www.sts.org/guidelines/lung-cancer-guidelines}
\end{enumerate}

\clearpage